\chapter{Graphs, trees and planarity}\label{ch:graphs}
\index{graph}

A graph is Chapter~4's relation made drawable: vertices, edges, and
nothing else --- yet from that skeleton come the oldest theorem of
network reasoning (handshake), the simplest connectivity story (trees),
Euler's walk through Königsberg, and a counting formula that forbids
whole graphs from living in the plane. After this chapter you can model
a problem as a graph, run degree-parity arguments, and use Euler's
formula as a non-planarity weapon.

\section{Graphs and the handshake lemma}\label{sec:handshake}
\index{degree}\index{handshake lemma}
\begin{definition}
A graph $G = (V, E)$ is a set of vertices with edges as unordered pairs
$\{u,v\}$; $\deg(v)$ counts the edges at $v$.
\end{definition}

\begin{example}\label{ex:handshake}
Prove that at any party, an even number of people shook an odd number
of hands.
\par\smallskip
\textbf{Strategy.} Count edge-ends two ways --- the double-counting
habit from Chapter~8: each edge contributes exactly $2$ to the degree
sum (Figure~\ref{fig:handshake}).
\par\smallskip
\begin{solution}
$\sum_v \deg(v) = 2|E|$ (each edge has two ends). An odd number of odd
terms can't sum to the even right side, so the number of odd-degree
vertices is even.
\end{solution}
\end{example}

\begin{figure}[ht]\centering
% Figure: handshake double counting -- each edge has two ends.
% Source for Figure~\ref{fig:handshake}. TikZ only; regenerable by compiling main.tex.
\begin{tikzpicture}[
  vx/.style={circle,draw=ink,fill=ink!10,inner sep=2.4pt,minimum size=8mm},
  lab/.style={font=\footnotesize\sffamily,text=slate},
  tick/.style={accent,font=\footnotesize\bfseries},
]
\node[vx] (a) at (0,1) {a};
\node[vx] (b) at (1.4,1.8) {b};
\node[vx] (c) at (2.8,1) {c};
\node[vx] (d) at (1.4,0) {d};
\draw[rulegray,thick] (a)--(b) (b)--(c) (a)--(d) (c)--(d) (b)--(d);
\node[tick] at (0.35,1.5) {+1}; \node[tick] at (0.95,1.62) {+1};
\node[tick] at (1.85,1.62) {+1}; \node[tick] at (2.45,1.5) {+1};
\node[tick] at (0.3,0.45) {+1}; \node[tick] at (1.1,0.28) {+1};
\node[tick] at (2.5,0.45) {+1}; \node[tick] at (1.7,0.28) {+1};
\node[tick] at (1.15,0.85) {+1}; \node[tick] at (1.65,0.85) {+1};
\node[lab] at (1.4,-0.8) {$\sum_v \deg(v) = 2\lvert E\rvert$: five edges, ten ends counted};
\end{tikzpicture}

\caption{The handshake lemma is a double count. Notice every edge
carries two $+1$ markers --- the sum of degrees is $2|E|$ because the
count is per endpoint, not per edge.}
\label{fig:handshake}
\end{figure}

\section{Walks, paths, cycles, connectivity}\label{sec:walks}
\index{walk}\index{path}\index{cycle}\index{connected}
A \emph{walk} allows repeats; a \emph{path} forbids repeated vertices; a
\emph{cycle} returns to its start (length $\ge 3$). A graph is
\emph{connected} when every pair is joined by a path; the maximal
connected pieces are its \emph{components}. The terms differ in exactly
what repetition is allowed --- a frequent exam-time collision.

\section{Bipartite graphs}\label{sec:bip}
\index{bipartite}
A graph is bipartite iff $V$ splits into two independent classes ---
equivalently, iff it contains \emph{no odd cycle} (an odd cycle needs a
third colour; a 2-colouring propagates consistently when no odd cycle
exists to break it).

\section{Trees}\label{sec:trees}
\index{tree}\index{spanning tree}
A tree is connected and acyclic. Every tree on $n \ge 2$ vertices has a
leaf (a degree-1 vertex): start at any vertex, walk without repeating.
The tree is acyclic, so the walk cannot revisit and must stop; the
endpoint has degree~$1$ --- an unvisited neighbor would extend the walk,
a visited one would close a cycle.

\begin{example}\label{ex:tree}
Prove a tree on $n$ vertices has exactly $n - 1$ edges.
\par\smallskip
\textbf{Strategy.} Remove the extremal object: a leaf exists, deleting
it keeps the graph a tree, and the count drops by exactly one ---
induction with the induction doing real work (Figure~\ref{fig:leaf}).
\par\smallskip
\begin{solution}
Induct on $n$; $n = 1$ gives $0$ edges. For $n \ge 2$ pick a leaf
$\ell$: $T - \ell$ is still connected (only $\ell$'s single neighbour
touched it) and acyclic, hence a tree on $n-1$ vertices, so by
induction it has $n - 2$ edges. Adding $\ell$'s edge back gives
$n - 1$.
\end{solution}
\end{example}

\begin{figure}[ht]\centering
% Figure: leaf-removal induction on a tree (3 steps).
% Source for Figure~\ref{fig:leaf}. TikZ only; regenerable by compiling main.tex.
\begin{tikzpicture}[
  vx/.style={circle,draw=ink,inner sep=2pt,minimum size=5.5mm},
  del/.style={circle,draw=accent,dashed,inner sep=2pt,minimum size=5.5mm},
  lab/.style={font=\footnotesize\sffamily,text=slate},
]
% step 1: tree on 5 vertices, leaf marked
\node[vx] (r1) at (0,0.8) {}; \node[vx] (a1) at (-0.6,0) {};
\node[vx] (b1) at (0.6,0) {}; \node[vx] (c1) at (-0.9,-0.8) {};
\node[del] (d1) at (-0.3,-0.8) {};
\draw[rulegray] (r1)--(a1) (r1)--(b1) (a1)--(c1) (a1)--(d1);
\node[lab] at (0,-1.5) {$n=5$: a leaf exists};
% step 2: after removal, tree on 4
\node[vx] (r2) at (3,0.8) {}; \node[vx] (a2) at (2.4,0) {};
\node[vx] (b2) at (3.6,0) {}; \node[vx] (c2) at (2.1,-0.8) {};
\draw[rulegray] (r2)--(a2) (r2)--(b2) (a2)--(c2);
\node[lab] at (3,-1.5) {$n-1=4$: still a tree};
% step 3: count edges
\node[lab] at (6.4,0) {$e_n = e_{n-1} + 1$};
\node[lab] at (6.4,-0.5) {$= (n-2) + 1 = n-1$};
\node[lab] at (6.4,-1.5) {edges shrink with vertices};
\end{tikzpicture}

\caption{Leaf removal as the induction step. Notice deleting a leaf
preserves treeness --- connectedness survives because a leaf's edge was
its only link.}
\label{fig:leaf}
\end{figure}

A spanning tree of a connected graph is a tree using all its vertices:
repeat edge-deletion on cycles until none remain, $n-1$ edges minimum.

\section{Eulerian trails}\label{sec:euler}
\index{Eulerian trail}
\begin{example}\label{ex:konig}
Decide whether Königsberg's seven bridges admit a walk using each
exactly once.
\par\smallskip
\textbf{Strategy.} Every \emph{intermediate} visit to a land mass
consumes two edges --- one in, one out --- so interior vertices need
even degree; only start and end may be odd.
\par\smallskip
\begin{solution}
The multigraph of Figure~\ref{fig:konig} has all four vertices odd.
An Eulerian trail allows exactly $0$ or $2$ odd vertices: no such walk
exists. The same parity test decides any graph (Listing~\ref{lst:euler}).
\end{solution}
\end{example}

\begin{figure}[ht]\centering
% Figure: Konigsberg schematic -> multigraph (4 land masses, 7 bridges).
% Source for Figure~\ref{fig:konig}. TikZ only; regenerable by compiling main.tex.
\begin{tikzpicture}[
  vx/.style={circle,draw=ink,fill=ink!10,inner sep=2.4pt,minimum size=9mm,font=\footnotesize},
  lab/.style={font=\footnotesize\sffamily,text=slate},
]
% schematic left
\draw[rulegray!50,fill=rulegray!12] (-3.2,-0.2) rectangle (-0.2,0.4); % top bank
\draw[rulegray!50,fill=rulegray!12] (-3.2,-2.2) rectangle (-0.2,-1.6); % bottom bank
\draw[rulegray!50,fill=rulegray!12] (-2.4,-0.6) ellipse (0.55 and 0.3); % island top
\draw[rulegray!50,fill=rulegray!12] (-1.0,-1.4) ellipse (0.55 and 0.3); % island bottom
\foreach \x/\y in {-2.4/-0.35,-2.4/-1.05,-1.0/-0.35,-1.0/-1.35,-2.4/-1.9,-1.0/-1.95,-1.7/-0.7}
  \draw[accent,thick] (\x,\y) -- ++(0.32,0.05);
\node[lab] at (-1.7,-2.9) {the map (schematic)};
% multigraph right
\node[vx] (A) at (1.4,0.8) {A};
\node[vx] (B) at (3.4,0.8) {B};
\node[vx] (C) at (1.4,-1.2) {C};
\node[vx] (D) at (3.4,-1.2) {D};
\draw[rulegray,thick] (A) to[bend left=25] (B);
\draw[rulegray,thick] (A) to[bend right=25] (B);
\draw[rulegray,thick] (A)--(C) (A)--(D) (B)--(C) (B)--(D) (C)--(D);
\node[lab] at (2.4,-2.2) {the multigraph: all four degrees odd};
\end{tikzpicture}

\caption{Königsberg as a multigraph. Notice the map's geometry is
irrelevant --- only ``which land masses share a bridge'' survives, which
is exactly the move that turns a puzzle into graph theory.}
\label{fig:konig}
\end{figure}

\begin{checkpoint}[Predict before you run]
Listing~\ref{lst:euler} checks connectivity and degree parity on
four graphs, printing one of three statuses. Predict which of \{Königsberg multigraph, 5-cycle, 4-path\}
gets an Eulerian circuit versus a mere trail.
\end{checkpoint}

\begin{lstlisting}[language=Python,caption={\texttt{euler\_check.py}:
DFS connectivity + odd-degree count (full script in
\texttt{code/euler\_check.py}).},label={lst:euler}]
odd = sum(len(adj[v]) % 2 for v in range(n))
\end{lstlisting}
Real output:
\begin{lstlisting}[numbers=none]
Konigsberg multigraph: no Eulerian trail
5-cycle: Eulerian circuit
path of 4: Eulerian trail (no circuit)
random 6v/8e: Eulerian trail (no circuit)
\end{lstlisting}

\section{Planar graphs and Euler's formula}\label{sec:planar}
\index{planar graph}\index{Euler's formula}
\begin{theorem}[Euler's formula]
A connected planar embedding obeys $V - E + F = 2$.
\end{theorem}
(Figure~\ref{fig:eulerf} works it twice.) Corollaries by double-counting
edge-face incidence: $E \le 3V - 6$ for simple planar graphs ($V \ge 3$),
and $E \le 2V - 4$ when the graph is triangle-free (every face needs
$\ge 4$ sides).

\begin{figure}[ht]\centering
% Figure: V-E+F=2 worked on two planar embeddings.
% Source for Figure~\ref{fig:eulerf}. TikZ only; regenerable by compiling main.tex.
\begin{tikzpicture}[
  vx/.style={circle,draw=ink,fill=ink!10,inner sep=2pt,minimum size=5mm},
  lab/.style={font=\footnotesize\sffamily,text=slate},
]
% embedding 1: square with diagonal
\node[vx] (a) at (0,1) {}; \node[vx] (b) at (1.4,1) {};
\node[vx] (c) at (1.4,-0.4) {}; \node[vx] (d) at (0,-0.4) {};
\draw[rulegray,thick] (a)--(b) (b)--(c) (c)--(d) (d)--(a) (a)--(c);
\node[lab] at (0.7,-0.9) {$V{=}4,\ E{=}5,\ F{=}3$};
% embedding 2: triangle with interior point
\node[vx] (p) at (3.6,1) {}; \node[vx] (q) at (5.4,1) {};
\node[vx] (rr) at (4.5,-0.6) {}; \node[vx] (s) at (4.5,0.35) {};
\draw[rulegray,thick] (p)--(q) (q)--(rr) (rr)--(p) (p)--(s) (q)--(s) (rr)--(s);
\node[lab] at (4.5,-0.9) {$V{=}4,\ E{=}6,\ F{=}4$};
\node[lab] at (2.6,-1.6) {$V - E + F = 4 - 5 + 3 = 2$ and $4 - 6 + 4 = 2$ ---
same formula, different drawings};
\end{tikzpicture}

\caption{Euler's formula checked on two embeddings. Notice $F$ counts
the outside face too --- the plane's ``outside'' is a face in every
sense the formula needs.}
\label{fig:eulerf}
\end{figure}

\begin{example}\label{ex:k33}
Prove $K_{3,3}$ is nonplanar.
\par\smallskip
\textbf{Strategy.} $K_{3,3}$ is bipartite, hence triangle-free: the
sharper bound $E \le 2V - 4$ applies. Count and contradict.
\par\smallskip
\begin{solution}
$V = 6$, $E = 9$; a bipartite planar graph needs $E \le 2\cdot 6 - 4
= 8 < 9$ --- impossible. (For $K_5$: $E = 10 > 3\cdot 5 - 6 = 9$; both
failures are in Figure~\ref{fig:k5k33}.)
\end{solution}
\end{example}

\begin{figure}[ht]\centering
% Figure: K5 and K_{3,3} drawn with their forced crossings marked.
% Source for Figure~\ref{fig:k5k33}. TikZ only; regenerable by compiling main.tex.
\begin{tikzpicture}[
  vx/.style={circle,draw=ink,fill=ink!10,inner sep=2pt,minimum size=4.5mm},
  cr/.style={accent,font=\footnotesize\bfseries},
  lab/.style={font=\footnotesize\sffamily,text=slate},
]
% K5 on pentagon
\foreach \i in {0,...,4}
  \node[vx] (k\i) at ({90+\i*72}:1.1) {};
\draw[rulegray!60] (k0)--(k1) (k0)--(k2) (k0)--(k3) (k0)--(k4)
                    (k1)--(k2) (k1)--(k3) (k1)--(k4)
                    (k2)--(k3) (k2)--(k4) (k3)--(k4);
\node[cr] at (0,0) {$\times$};
\node[lab] at (0,-1.7) {$K_5$: $e = 10 > 3v - 6 = 9$};
% K33 two columns
\foreach \i in {0,1,2}
{
  \node[vx] (l\i) at (3.2,{0.9-\i*0.9}) {};
  \node[vx] (r\i) at (4.8,{0.9-\i*0.9}) {};
}
\foreach \i in {0,1,2}
  \foreach \j in {0,1,2}
    \draw[rulegray!60] (l\i)--(r\j);
\node[cr] at (4.0,0) {$\times$};
\node[lab] at (4.0,-1.7) {$K_{3,3}$: $e = 9 > 2v - 4 = 8$};
\node[lab] at (2.4,-2.4) {crossings drawn are heuristic --- the proofs are the edge-count contradictions};
\end{tikzpicture}

\caption{$K_5$ and $K_{3,3}$ with their inevitable crossings marked.
Notice the proofs are edge-counts, not the drawing --- the picture is
intuition, Euler's formula is the licence.}
\label{fig:k5k33}
\end{figure}

\section{Graph colouring basics}\label{sec:color}
\index{graph colouring}
A proper colouring assigns colours so no edge is monochromatic;
$\chi(G)$ is the fewest colours needed. Bipartite $\Rightarrow$
$\chi = 2$ (the definition repackaged); odd cycles need $3$; and a
vertex of degree $d$ bounds $\chi \le d_{\max} + 1$ greedily.

\begin{exercise}[easy]\label{exr:g1}
A graph has degrees $3,3,2,2,1,1$. Is it possible? Check the handshake
lemma first.
\end{exercise}
\begin{exercise}[medium]\label{exr:g2}
Prove a tree is bipartite.
\end{exercise}
\begin{exercise}[medium]\label{exr:g3}
Does the graph $K_4$ minus one edge have an Eulerian trail? A circuit?
Decide by parity, don't draw the trail.
\end{exercise}
\begin{exercise}[medium]\label{exr:g4}
Use Euler's formula to show every planar simple graph has a vertex of
degree $\le 5$.
\end{exercise}

\section{Chapter summary}
\begin{table}[ht]\centering
\begin{tabularx}{\linewidth}{@{}lYY@{}}
\toprule
\tablehead\thead{Fact} & \thead{Statement} & \thead{Proof tool}\\
\midrule
handshake & $\sum\deg = 2|E|$ & double counting\\
tree & $n$ vertices $\Rightarrow$ $n-1$ edges & leaf induction\\
bipartite & $\Leftrightarrow$ no odd cycle & 2-colour propagation\\
Euler trail & conn.\ + $\le 2$ odd degrees & parity argument\\
Euler formula & $V-E+F=2$ planar & induction on edges\\
planar bound & $E\le 3V-6$; $2V-4$ if no $\triangle$ & face-edge count\\
\bottomrule
\end{tabularx}
\caption{One page: degrees and double counts, then Euler twice ---
trails by parity, planarity by formula.}
\label{tab:graphs}
\end{table}

\section{Common mistakes}
\index{common mistakes}
\begin{itemize}
\item \textbf{Walk $\ne$ path $\ne$ cycle.} A walk may repeat; a path
may not; a cycle returns home. Many ``proofs'' break by silently
upgrading a walk to a path.
\item \textbf{Using $E \le 3V-6$ backwards.} Satisfying the bound does
\emph{not} prove planarity --- it is a necessary condition only;
nonplanarity is proved by violating it.
\item \textbf{``Looks bipartite.''} The property is odd-cycle-free, not
visually two-sided; a twisted drawing can hide an odd cycle.
\end{itemize}

\section{Additional exercises}
\begin{exercise}[easy]
Can a graph have degree sequence $5,4,3,2,1,0$? What invariant decides
before you try to draw it?
\end{exercise}
\begin{exercise}[easy]
How many edges does a forest of $k$ components on $n$ vertices have?
Generalise the tree count.
\end{exercise}
\begin{exercise}[easy]
Is the cube graph (8 vertices, edges of a cube) bipartite? Eulerian?
\end{exercise}
\begin{exercise}[easy]
How many spanning trees does a cycle $C_n$ have?
\end{exercise}
\begin{exercise}[medium]
Prove every connected graph with all even degrees has an Eulerian
circuit (Hierholzer's idea: stitch cycles).
\end{exercise}
\begin{exercise}[medium]
Show a tree with a vertex of degree $d$ has at least $d$ leaves.
\end{exercise}
\begin{exercise}[medium]
Prove a connected graph where every vertex has degree $2$ is a cycle.
\end{exercise}
\begin{exercise}[medium]
The Petersen graph has $V=10$, $E=15$. Is it planar? (Its girth is 5 ---
use the right bound.)
\end{exercise}
\begin{exercise}[hard]
Prove $K_5$ nonplanar without citing the $3V-6$ bound --- derive the
bound first from $V-E+F=2$.
\end{exercise}
\begin{exercise}[hard]
Show that a planar graph and its dual have complementary edge counts,
and that the dodecahedron's dual is the icosahedron (20 faces, 12
vertices).
\end{exercise}
\begin{exercise}[hard]
Extend \texttt{euler\_check.py} to output one actual Eulerian trail
(Hierholzer's algorithm) and verify it uses each edge exactly once on
200 random even-degree graphs.
\end{exercise}

\chapternotes
The walk/path/cycle hygiene and the planar sections follow 6.042J's
graph chapters~\cite{mcs}; CS70's graph notes carry the same material
with more algorithms~\cite{cs70}; Engel's graph chapter supplies the
olympiad drill~\cite{engel}, and Cut-the-Knot the Königsberg
framing~\cite{cuttheknot}. Contest items powered by this chapter ---
graph-distance and path-state framings (AMC 12 2023 A P21,
AIME 2024 I P3), cycle and polygon walks (AIME 2018 I P10,
AIME 2018 I P14, AIME 2021 I P12), plus constraint graphs
(AMC 12 2024 A P22, AIME 2023 I P3) --- are on the AoPS
archive~\cite{aops-amc}.
